\part{紧凑动态风格}

\chapter{数据结构}

\Needspace{280pt}

\section{并查集}\label{compact-DSU}\index{DSU}

init(n)、find(x)、merge(x,y)，点号 1..n；siz[find(x)] 为集合大小。路径压缩与按大小合并，均摊 O($\alpha$(n))。

\lstinputlisting[firstline=6,lastline=28]{../src/compact/data_structure.hpp}

\Needspace{323pt}

\section{可撤销并查集}\label{compact-RollbackDSU}\index{RollbackDSU}

构造后 snapshot() 保存历史长度，rollback(t) 回退。禁止路径压缩；find/merge 为 O(log n)，回退每次成功合并 O(1)。

\lstinputlisting[firstline=29,lastline=56]{../src/compact/data_structure.hpp}

\Needspace{255pt}

\section{树状数组与第 k 小}\label{compact-Fenwick}\index{Fenwick}

add(x,v)、sum(x)、query(l,r)，下标 1..n，闭区间；kth(k) 要求所有频率非负且 k>0，超过总和返回 n+1。单次 O(log n)。

\lstinputlisting[firstline=57,lastline=76]{../src/compact/data_structure.hpp}

\Needspace{358pt}

\section{区间加与区间和}\label{compact-LazySeg}\index{LazySeg}

初值全零；add(l,r,v)、query(l,r)，1$\le$l$\le$r$\le$n。维护 sum 与加法 tag；所有中间和与长度乘积须在 long long 范围内。单次 O(log n)。

\lstinputlisting[firstline=77,lastline=108]{../src/compact/data_structure.hpp}

\Needspace{210pt}

\section{64 位异或线性基}\label{compact-XorBasis}\index{XorBasis}

insert(x) 返回是否增加秩；query(x) 返回 x 与所张成空间异或后的最大无符号值。单次 O(64)，不提供第 k 小。

\lstinputlisting[firstline=109,lastline=123]{../src/compact/data_structure.hpp}

\chapter{网络流}

\section{最大流与方案}\label{compact-Dinic}\index{Dinic}

Dinic(n)，add(u,v,c) 返回偶数原边编号；flow(s,t,limit) 返回本次新增流；used(id) 查询原边净流，cut(s) 返回残量可达点集。允许重边、自环，容量非负，s$\ne$t；重复求流须使用相同源汇。只有求尽最大流后 cut 才是最小割。一般 O(V$^{2}$E)，递归增广深度 O(V)。总流不得超过 long long 上限。

\lstinputlisting[firstline=6,lastline=67]{../src/compact/flow.hpp}

\section{势能最小费用最大流}\label{compact-MinCostFlow}\index{MinCostFlow}

add(u,v,cap,cost) 返回边编号，flow(s,t,limit) 返回新增流量及 128 位费用，used(id) 返回原边流量。允许负费用，初始残量网络必须没有负环；检测到负环抛异常。Bellman-Ford 初始化 O(VE)，每次增广 O(E log V)；增广次数不能简单假设为 O(V)。 另要求 $\sum$(cap×|cost|)<2\textasciicircum{}120，防止总费用超过 int128。

\lstinputlisting[firstline=68,lastline=131]{../src/compact/flow.hpp}

\Needspace{286pt}

\section{上下界可行循环流}\label{compact-BoundedCirculation}\index{BoundedCirculation}

add(u,v,lo,hi)，0$\le$lo$\le$hi；solve() 仅调用一次，返回可行性，仅成功后读取 used(id)。约束为每点流量守恒；不是带源汇最大流。所有上下界及其总和须在 long long 范围内。依赖 Dinic。

\lstinputlisting[firstline=132,lastline=154]{../src/compact/flow.hpp}

\chapter{图论}

\Needspace{344pt}

\section{非负权最短路}\label{compact-Dijkstra}\index{Dijkstra}

add(u,v,w)、run(s)，dis 为距离，pre 为前驱，path(t) 返回点序列。不可达用 LLONG\_MAX；所有可表示的最短距离必须严格小于该哨兵，w$\ge$0。O((V+E)log V)。

\lstinputlisting[firstline=6,lastline=35]{../src/compact/graph.hpp}

\Needspace{400pt}

\section{强连通分量}\label{compact-SCC}\index{SCC}

add(u,v)、run()，bel[1..n] 为分量编号，cnt 为数量；分量编号按缩点图拓扑序递增。迭代 Kosaraju 避免递归爆栈；O(V+E)，同时存正反图。

\lstinputlisting[firstline=36,lastline=72]{../src/compact/graph.hpp}

\Needspace{248pt}

\section{2-SAT 与任意可行方案}\label{compact-TwoSAT}\index{TwoSAT}

add(x,a,y,b) 添加 (x=a) 或 (y=b)，变量 1..n；solve() 返回可满足性，仅成功时读取 ans。依赖同文件 SCC；O(V+E)，不保证字典序最小。

\lstinputlisting[firstline=73,lastline=91]{../src/compact/graph.hpp}

\Needspace{381pt}

\section{二分图最大匹配}\label{compact-BipartiteMatching}\index{BipartiteMatching}

左右两侧分别编号 1..n、1..m；add(u,v)，solve() 返回匹配数，l[u] 与 r[v] 为匹配端点，0 表示未匹配。分层增广；本版本保守按 O(VE) 估计，不宣称已经验证 Hopcroft-Karp 最紧复杂度。

\lstinputlisting[firstline=92,lastline=125]{../src/compact/graph.hpp}

\Needspace{334pt}

\section{割点与桥}\label{compact-Lowlink}\index{Lowlink}

add(u,v) 返回无向边编号；run() 后 cut[u] 标识割点、bridge[id] 标识桥。按边编号跳过父边，支持重边、自环、不连通图。O(V+E)，递归深度可能 O(V)。

\lstinputlisting[firstline=126,lastline=154]{../src/compact/graph.hpp}

\chapter{树上算法}

\Needspace{555pt}

\section{重链剖分与 LCA}\label{compact-HLD}\index{HLD}

add(u,v) 建无向树；build(root) 后读取 fa、dep、siz、son、top、dfn、rk。子树对应 [dfn[u],dfn[u]+siz[u]-1]。lca(u,v) 为 O(log n)，path(u,v,work,edge) 分解路径；work(l,r) 仅适用于可交换聚合，边权模式排除 LCA 位置。输入必须是连通无环树，建树 O(n)，无递归。

\lstinputlisting[firstline=6,lastline=58]{../src/compact/tree.hpp}

\chapter{字符串}

\section{KMP、Z、Manacher 与最小表示}\label{compact-StringAlgo}\index{StringAlgo}

字符串及返回位置均 0-based。prefix(s)、z(s)；match(s,t) 查非空模式串并返回重叠匹配起点。manacher 返回奇数半径（含中心）和偶数半径（中心在 i 前）；rotation 返回最小循环表示起点，空串返回 0。各 O(n+m)。

\lstinputlisting[firstline=6,lastline=77]{../src/compact/string.hpp}

\Needspace{424pt}

\section{AC 自动机出现次数}\label{compact-AhoCorasick}\index{AhoCorasick}

字母表 a..z，add(s) 返回终止状态，要求模式非空；全部插入后 build()，只允许一次构建。count(text) 返回每个状态对应模式的出现次数，可重复查询；重复模式共享状态。构建 O(26N)，查询 O(|text|+N)。

\lstinputlisting[firstline=78,lastline=116]{../src/compact/string.hpp}

\Needspace{320pt}

\section{倍增后缀数组与 LCP}\label{compact-SuffixArray}\index{SuffixArray}

sa[k] 为第 k 个后缀起点，rk[i] 为排名，lcp[k] 为 sa[k-1] 与 sa[k] 的 LCP（lcp[0]=0）。空串合法；当前排序比较版 O(n log$^{2}$ n)，不能替代待补的线性 SA/DC3 条目。

\lstinputlisting[firstline=117,lastline=143]{../src/compact/string.hpp}

\Needspace{432pt}

\section{后缀自动机}\label{compact-SuffixAutomaton}\index{SuffixAutomaton}

extend(c) 接收整数 0..25；distinct() 返回不同非空子串数，counts() 返回 endpos 次数且不修改结构，可重复调用。状态 a，初始状态 0，link 根为 -1；空间 O(26n)，构建 O(26n)。

\lstinputlisting[firstline=144,lastline=183]{../src/compact/string.hpp}

\chapter{数论}

\section{模运算、素性、CRT 与类欧几里德}\label{compact-NumberTheory}\index{NumberTheory}

mul/power 支持无符号 64 位正模数，乘积用 uint128。prime 使用确定性 64 位底数。inverse(a,m) 无逆元返回 -1。crt(r,m,b,n) 合并同余式，无解返回 false，模数溢出抛异常。floor\_sum 为 i=0..n-1 的 floor((ai+b)/m) 之和，支持负 a、b，返回 int128；约定 n$\le$10$^{9}$、m$\le$10$^{9}$，a、b 为有符号 64 位。

\lstinputlisting[firstline=6,lastline=74]{../src/compact/number_theory.hpp}

\Needspace{327pt}

\section{Pollard–Rho 因数分解}\label{compact-PollardRho}\index{PollardRho}

factor(n) 返回升序质因子（含重数），n$\ge$1；依赖 NumberTheory。返回因子经过整除和素性检查；Las Vegas 随机运行时间，固定默认种子方便复现，可由调用者更换。无确定性最坏时间保证。

\lstinputlisting[firstline=75,lastline=102]{../src/compact/number_theory.hpp}

\Needspace{252pt}

\section{线性筛、欧拉函数与 Möbius 函数}\label{compact-LinearSieve}\index{LinearSieve}

构造参数 n$\ge$0，prime 为质数表，lp/phi/mu 为 0..n 数组。phi[1]=mu[1]=1，O(n) 时间和空间。

\lstinputlisting[firstline=103,lastline=122]{../src/compact/number_theory.hpp}

\Needspace{196pt}

\section{静态模整数}\label{compact-ModInt}\index{ModInt}

模板参数必须为不超过 INT\_MAX 的素数且$\ge$2。加减乘、pow、inv、除法；求逆要求非零，pow 指数非负。复杂度：加减乘 O(1)，幂和逆 O(log mod)。

\lstinputlisting[firstline=123,lastline=135]{../src/compact/number_theory.hpp}

\Needspace{199pt}

\section{阶乘组合数}\label{compact-Binomial}\index{Binomial}

Binomial<prime>(n)，要求 0$\le$n<prime；choose(n,k) 越界 k 返回 0。预处理 O(n+log prime)，单次 O(1)。尚不能替代合数模数 exLucas。

\lstinputlisting[firstline=136,lastline=148]{../src/compact/number_theory.hpp}

\chapter{多项式}

\section{NTT、形式幂级数、FWT 与递推}\label{compact-Polynomial}\index{Polynomial}

依赖 number\_theory.hpp。模数固定 998244353，乘法结果长度$\le$2$^{23}$；inverse 要求 a[0]$\ne$0，log 要求 a[0]=1，exp 要求 a[0]=0，截断长度 n>0。朴素逐层 NTT 实现的 FPS 为 O(n log n)。FWT 的 op 为 \textasciicircum{}、|、\&；BM 返回 s[n]=$\sum$c[i-1]s[n-i] 的系数；recurrence 用 O(k$^{2}$ log n) 求第 n 项。BM 对有限样本给出的递推不自动保证未知后续项成立。

\lstinputlisting[firstline=5,lastline=129]{../src/compact/polynomial.hpp}

\chapter{代数与数论进阶}

\section{模高斯消元、行列式、矩阵树定理}\label{compact-LinearAlgebra}\index{LinearAlgebra}

素数模数。solve(a,n) 输入增广矩阵，返回 consistent、rank、特解 particular 和零空间基 kernel。determinant(empty)=1。spanning\_trees(n,edges) 使用 0-based 无向图，允许重边，忽略自环；返回生成树数模 p。高斯与行列式 O(n$^{3}$)，矩阵幂 O(n$^{3}$ log e)。

\lstinputlisting[firstline=4,lastline=79]{../src/compact/algebra.hpp}

\Needspace{329pt}

\section{杜教筛}\label{compact-DuJiao}\index{DuJiao}

构造预筛上限 B$\ge$1，mertens(n) 返回 $\mu$ 的前缀和，totient\_sum(n) 返回 $\varphi$ 的前缀和（int128）。约定 0$\le$n$\le$10$^{11}$，按整除分块递归并缓存；大输入需要合理 B，不能把常数 B 当作性能保证。

\lstinputlisting[firstline=80,lastline=107]{../src/compact/algebra.hpp}

\Needspace{325pt}

\section{扩展 BSGS}\label{compact-DiscreteLog}\index{DiscreteLog}

solve(a,b,m) 返回最小 x$\ge$0 使 a\textasciicircum{}x$\equiv$b (mod m)，无解 -1；1$\le$m$\le$10$^{12}$，0\textasciicircum{}0=1。时间与内存 O($\sqrt{\vphantom x}$m)，unordered\_map 复杂度为期望界。

\lstinputlisting[firstline=108,lastline=135]{../src/compact/algebra.hpp}

\chapter{计算几何}

\section{整数精确二维几何}\label{compact-IntegerGeometry}\index{IntegerGeometry}

坐标绝对值$\le$10$^{12}$，点数$\le$10$^{6}$。叉积、点积、线段相交、凸包、点在多边形、点在凸包、旋转卡壳直径。凸包逆时针，不保留共线内点，不重复首点。contains 返回 0 外部、1 边界、2 内部。convex\_contains/diameter2 输入必须为本模板 strict hull。PolarLess 要求排除零向量。

\lstinputlisting[firstline=6,lastline=109]{../src/compact/geometry.hpp}

\section{浮点点线圆构造}\label{compact-RealGeometry}\index{RealGeometry}

long double，显式 eps 默认 10$^{-12}$。要求有限输入及适当尺度；接近退化时按容差分类，不能提供任意实数输入的精确拓扑保证。Result 区分 none/one/two/infinite/degenerate。零长度线返回 degenerate；投影与点到线段距离将零线段视为点。坐标缩放和误差预算应按具体题面复核。

\lstinputlisting[firstline=110,lastline=189]{../src/compact/geometry.hpp}

\chapter{优化与可持久化}

\Needspace{287pt}

\section{Max-plus 矩阵幂}\label{compact-MaxPlusMatrix}\index{MaxPlusMatrix}

a[i][j] 为转移收益，不存在的边为 neg；power(k) 求恰好 k 步的最优转移，单位矩阵对角为 0。所有有限中间值绝对值小于 2\textasciicircum{}119。O(n$^{3}$ log k)；2025 成都 B 的转移类型，具体题目的建模仍需调用者完成。

\lstinputlisting[firstline=6,lastline=28]{../src/compact/optimization.hpp}

\Needspace{492pt}

\section{离散李超树与最优直线编号}\label{compact-LiChao}\index{LiChao}

构造传入所有可能查询的横坐标，add(\{k,b,id\}) 插入直线，query(x) 返回最小值与最小编号，要求 x 在预登记坐标内且至少有一条线。整数求值用 int128，k、b、x 为有符号 64 位且中间值可表示。单次 O(log n)。

\lstinputlisting[firstline=29,lastline=75]{../src/compact/optimization.hpp}

\Needspace{386pt}

\section{主席树：静态区间第 k 小}\label{compact-PersistentKth}\index{PersistentKth}

构造传入 0-based 数组，kth(l,r,k) 的区间使用 1-based 闭区间；要求 1$\le$k$\le$r-l+1。离散化后每个前缀一个版本，单次 O(log n)，总空间 O(n log n)。

\lstinputlisting[firstline=76,lastline=109]{../src/compact/optimization.hpp}

\chapter{图论进阶}

\Needspace{389pt}

\section{矩形二分图最小权匹配}\label{compact-Hungarian}\index{Hungarian}

solve(cost) 输入 n×m 完整权值矩阵，0$\le$n$\le$m；每行匹配不同列，返回最小费用及列号（0-based）。允许负费用，使用 int128 势能；O(n$^{2}$m)。禁止边不能用未经证明的有限大常数随意替代，当前接口仅表示完整二分图。

\lstinputlisting[firstline=6,lastline=40]{../src/compact/graph_advanced.hpp}

\Needspace{390pt}

\section{有向最小树形图}\label{compact-Arborescence}\index{Arborescence}

solve(n,root,edges)，点号 0..n-1，n$\ge$1；返回以 root 为根向外可达的最小树形图费用，无解 nullopt。允许负权、重边、自环（自环不选）；O(VE)。当前仅返回费用，不恢复边方案。所有中间值须远小于 2\textasciicircum{}120。

\lstinputlisting[firstline=41,lastline=75]{../src/compact/graph_advanced.hpp}

\Needspace{340pt}

\section{无向全局最小割}\label{compact-StoerWagner}\index{StoerWagner}

solve(w)，n$\ge$2；w 是对称非负权矩阵，重边需先累加，主对角置零。返回割权与一侧顶点编号（0-based），允许不连通图。O(n$^{3}$)，总权值小于 2\textasciicircum{}120。

\lstinputlisting[firstline=76,lastline=105]{../src/compact/graph_advanced.hpp}

\chapter{精确几何进阶}

\section{整数最近点对与 Minkowski 和}\label{compact-GeometryExtra}\index{GeometryExtra}

closest\_pair(points) 返回最短距离平方，少于两个点为 nullopt，重复点返回 0，O(n log n)，全程精确整数。minkowski 输入 strict CCW 凸包，返回不重复首点的凸包；非退化 O(n+m)，退化时 O((n+m)log(n+m))。坐标及坐标和须满足整数几何范围。

\lstinputlisting[firstline=4,lastline=68]{../src/compact/geometry_extra.hpp}

\Needspace{233pt}

\section{整数三维共线、共面与方向}\label{compact-IntegerGeometry3D}\index{IntegerGeometry3D}

坐标绝对值$\le$10$^{9}$；orient(a,b,c,d) 为有向六倍四面体体积，0 表示共面；collinear 判三点共线，on\_segment 支持退化线段。int128 精确判定，不包含三维凸包构造。

\lstinputlisting[firstline=69,lastline=85]{../src/compact/geometry_extra.hpp}

\part{传统接口风格}

\chapter{数据结构}

\Needspace{365pt}

\section{并查集}\label{classic-DSU}\index{Disjoint\_Set}

算法约束见紧凑版本第~\pageref{compact-DSU}~页。本节代码独立可抄；采用下列实际接口名。

先声明容量模板参数，再调用 Init(n) 初始化；大对象必须置于全局或 static 存储。

\lstinputlisting[firstline=10,lastline=42]{../src/classic/data_structure.hpp}

\Needspace{382pt}

\section{可撤销并查集}\label{classic-RollbackDSU}\index{Rollback\_DSU}

算法约束见紧凑版本第~\pageref{compact-RollbackDSU}~页。本节代码独立可抄；采用下列实际接口名。

\lstinputlisting[firstline=123,lastline=157]{../src/classic/data_structure.hpp}

\Needspace{374pt}

\section{树状数组与第 k 小}\label{classic-Fenwick}\index{Binary\_Indexed\_Tree}

算法约束见紧凑版本第~\pageref{compact-Fenwick}~页。本节代码独立可抄；采用下列实际接口名。

先声明容量模板参数，再调用 Init(n) 初始化；大对象必须置于全局或 static 存储。

\lstinputlisting[firstline=43,lastline=76]{../src/classic/data_structure.hpp}

\Needspace{477pt}

\section{区间加与区间和}\label{classic-LazySeg}\index{Segment\_Tree}

算法约束见紧凑版本第~\pageref{compact-LazySeg}~页。本节代码独立可抄；采用下列实际接口名。

先声明容量模板参数，再调用 Init(n) 初始化；大对象必须置于全局或 static 存储。

\lstinputlisting[firstline=77,lastline=122]{../src/classic/data_structure.hpp}

\Needspace{244pt}

\section{64 位异或线性基}\label{classic-XorBasis}\index{Xor\_Basis}

算法约束见紧凑版本第~\pageref{compact-XorBasis}~页。本节代码独立可抄；采用下列实际接口名。

\lstinputlisting[firstline=158,lastline=176]{../src/classic/data_structure.hpp}

\chapter{网络流}

\section{最大流与方案}\label{classic-Dinic}\index{Network\_Flow}

算法约束见紧凑版本第~\pageref{compact-Dinic}~页。本节代码独立可抄；采用下列实际接口名。

先声明容量模板参数，再调用 Init(n) 初始化；大对象必须置于全局或 static 存储。

\lstinputlisting[firstline=9,lastline=79]{../src/classic/flow.hpp}

\section{势能最小费用最大流}\label{classic-MinCostFlow}\index{Min\_Cost\_Flow}

算法约束见紧凑版本第~\pageref{compact-MinCostFlow}~页。本节代码独立可抄；采用下列实际接口名。

使用 Insert 加边、Flow 求流、Used 查看方案。

\lstinputlisting[firstline=21,lastline=117]{../src/classic/min_cost_flow.hpp}

\Needspace{414pt}

\section{上下界可行循环流}\label{classic-BoundedCirculation}\index{Bounded\_Circulation}

算法约束见紧凑版本第~\pageref{compact-BoundedCirculation}~页。本节代码独立可抄；采用下列实际接口名。

\lstinputlisting[firstline=80,lastline=117]{../src/classic/flow.hpp}

\chapter{图论}

\Needspace{497pt}

\section{非负权最短路}\label{classic-Dijkstra}\index{Shortest\_Path}

算法约束见紧凑版本第~\pageref{compact-Dijkstra}~页。本节代码独立可抄；采用下列实际接口名。

\lstinputlisting[firstline=22,lastline=69]{../src/classic/graph.hpp}

\section{强连通分量}\label{classic-SCC}\index{Strong\_Component}

算法约束见紧凑版本第~\pageref{compact-SCC}~页。本节代码独立可抄；采用下列实际接口名。

\lstinputlisting[firstline=70,lastline=132]{../src/classic/graph.hpp}

\Needspace{333pt}

\section{2-SAT 与任意可行方案}\label{classic-TwoSAT}\index{Two\_SAT}

算法约束见紧凑版本第~\pageref{compact-TwoSAT}~页。本节代码独立可抄；采用下列实际接口名。

\lstinputlisting[firstline=133,lastline=161]{../src/classic/graph.hpp}

\section{二分图最大匹配}\label{classic-BipartiteMatching}\index{Bipartite\_Matching}

算法约束见紧凑版本第~\pageref{compact-BipartiteMatching}~页。本节代码独立可抄；采用下列实际接口名。

\lstinputlisting[firstline=162,lastline=225]{../src/classic/graph.hpp}

\Needspace{504pt}

\section{割点与桥}\label{classic-Lowlink}\index{Low\_Link}

算法约束见紧凑版本第~\pageref{compact-Lowlink}~页。本节代码独立可抄；采用下列实际接口名。

\lstinputlisting[firstline=226,lastline=274]{../src/classic/graph.hpp}

\chapter{树上算法}

\section{重链剖分与 LCA}\label{classic-HLD}\index{Heavy\_Light\_Decomposition}

算法约束见紧凑版本第~\pageref{compact-HLD}~页。本节代码独立可抄；采用下列实际接口名。

\lstinputlisting[firstline=22,lastline=111]{../src/classic/tree.hpp}

\chapter{字符串}

\section{KMP、Z、Manacher 与最小表示}\label{classic-StringAlgo}\index{String\_Algorithm}

算法约束见紧凑版本第~\pageref{compact-StringAlgo}~页。本节代码独立可抄；采用下列实际接口名。

\lstinputlisting[firstline=22,lastline=136]{../src/classic/string.hpp}

\section{AC 自动机出现次数}\label{classic-AhoCorasick}\index{AC\_Automaton}

算法约束见紧凑版本第~\pageref{compact-AhoCorasick}~页。本节代码独立可抄；采用下列实际接口名。

\lstinputlisting[firstline=137,lastline=206]{../src/classic/string.hpp}

\Needspace{473pt}

\section{倍增后缀数组与 LCP}\label{classic-SuffixArray}\index{Suffix\_Array}

算法约束见紧凑版本第~\pageref{compact-SuffixArray}~页。本节代码独立可抄；采用下列实际接口名。

\lstinputlisting[firstline=207,lastline=251]{../src/classic/string.hpp}

\section{后缀自动机}\label{classic-SuffixAutomaton}\index{Suffix\_Automaton}

算法约束见紧凑版本第~\pageref{compact-SuffixAutomaton}~页。本节代码独立可抄；采用下列实际接口名。

\lstinputlisting[firstline=252,lastline=322]{../src/classic/string.hpp}

\chapter{数论}

\section{模运算、素性、CRT 与类欧几里德}\label{classic-NumberTheory}\index{Number\_Theory}

算法约束见紧凑版本第~\pageref{compact-NumberTheory}~页。本节代码独立可抄；采用下列实际接口名。

\lstinputlisting[firstline=22,lastline=155]{../src/classic/number_theory.hpp}

\Needspace{497pt}

\section{Pollard–Rho 因数分解}\label{classic-PollardRho}\index{Pollard\_Rho}

算法约束见紧凑版本第~\pageref{compact-PollardRho}~页。本节代码独立可抄；采用下列实际接口名。

\lstinputlisting[firstline=156,lastline=203]{../src/classic/number_theory.hpp}

\Needspace{380pt}

\section{线性筛、欧拉函数与 Möbius 函数}\label{classic-LinearSieve}\index{Linear\_Sieve}

算法约束见紧凑版本第~\pageref{compact-LinearSieve}~页。本节代码独立可抄；采用下列实际接口名。

\lstinputlisting[firstline=204,lastline=238]{../src/classic/number_theory.hpp}

\Needspace{392pt}

\section{静态模整数}\label{classic-ModInt}\index{Mod\_Int}

算法约束见紧凑版本第~\pageref{compact-ModInt}~页。本节代码独立可抄；采用下列实际接口名。

\lstinputlisting[firstline=239,lastline=274]{../src/classic/number_theory.hpp}

\Needspace{275pt}

\section{阶乘组合数}\label{classic-Binomial}\index{Combination}

算法约束见紧凑版本第~\pageref{compact-Binomial}~页。本节代码独立可抄；采用下列实际接口名。

\lstinputlisting[firstline=275,lastline=296]{../src/classic/number_theory.hpp}

\chapter{多项式}

\section{NTT、形式幂级数、FWT 与递推}\label{classic-Polynomial}\index{Polynomial}

算法约束见紧凑版本第~\pageref{compact-Polynomial}~页。本节代码独立可抄；采用下列实际接口名。

\lstinputlisting[firstline=5,lastline=224]{../src/classic/polynomial.hpp}

\chapter{代数与数论进阶}

\section{模高斯消元、行列式、矩阵树定理}\label{classic-LinearAlgebra}\index{Linear\_Algebra}

算法约束见紧凑版本第~\pageref{compact-LinearAlgebra}~页。本节代码独立可抄；采用下列实际接口名。

\lstinputlisting[firstline=4,lastline=132]{../src/classic/algebra.hpp}

\Needspace{499pt}

\section{杜教筛}\label{classic-DuJiao}\index{Du\_Jiao}

算法约束见紧凑版本第~\pageref{compact-DuJiao}~页。本节代码独立可抄；采用下列实际接口名。

\lstinputlisting[firstline=133,lastline=180]{../src/classic/algebra.hpp}

\Needspace{470pt}

\section{扩展 BSGS}\label{classic-DiscreteLog}\index{Discrete\_Log}

算法约束见紧凑版本第~\pageref{compact-DiscreteLog}~页。本节代码独立可抄；采用下列实际接口名。

\lstinputlisting[firstline=181,lastline=225]{../src/classic/algebra.hpp}

\chapter{计算几何}

\section{整数精确二维几何}\label{classic-IntegerGeometry}\index{Integer\_Geometry}

算法约束见紧凑版本第~\pageref{compact-IntegerGeometry}~页。本节代码独立可抄；采用下列实际接口名。

\lstinputlisting[firstline=22,lastline=183]{../src/classic/geometry.hpp}

\section{浮点点线圆构造}\label{classic-RealGeometry}\index{Real\_Geometry}

算法约束见紧凑版本第~\pageref{compact-RealGeometry}~页。本节代码独立可抄；采用下列实际接口名。

\lstinputlisting[firstline=184,lastline=332]{../src/classic/geometry.hpp}

\chapter{优化与可持久化}

\Needspace{381pt}

\section{Max-plus 矩阵幂}\label{classic-MaxPlusMatrix}\index{Max\_Plus\_Matrix}

算法约束见紧凑版本第~\pageref{compact-MaxPlusMatrix}~页。本节代码独立可抄；采用下列实际接口名。

\lstinputlisting[firstline=22,lastline=55]{../src/classic/optimization.hpp}

\section{离散李超树与最优直线编号}\label{classic-LiChao}\index{Li\_Chao\_Tree}

算法约束见紧凑版本第~\pageref{compact-LiChao}~页。本节代码独立可抄；采用下列实际接口名。

\lstinputlisting[firstline=56,lastline=146]{../src/classic/optimization.hpp}

\section{主席树：静态区间第 k 小}\label{classic-PersistentKth}\index{Persistent\_Kth}

算法约束见紧凑版本第~\pageref{compact-PersistentKth}~页。本节代码独立可抄；采用下列实际接口名。

\lstinputlisting[firstline=147,lastline=211]{../src/classic/optimization.hpp}

\chapter{图论进阶}

\section{矩形二分图最小权匹配}\label{classic-Hungarian}\index{Kuhn\_Munkres}

算法约束见紧凑版本第~\pageref{compact-Hungarian}~页。本节代码独立可抄；采用下列实际接口名。

\lstinputlisting[firstline=22,lastline=87]{../src/classic/graph_advanced.hpp}

\section{有向最小树形图}\label{classic-Arborescence}\index{Directed\_MST}

算法约束见紧凑版本第~\pageref{compact-Arborescence}~页。本节代码独立可抄；采用下列实际接口名。

\lstinputlisting[firstline=88,lastline=153]{../src/classic/graph_advanced.hpp}

\Needspace{519pt}

\section{无向全局最小割}\label{classic-StoerWagner}\index{Global\_Min\_Cut}

算法约束见紧凑版本第~\pageref{compact-StoerWagner}~页。本节代码独立可抄；采用下列实际接口名。

\lstinputlisting[firstline=154,lastline=204]{../src/classic/graph_advanced.hpp}

\chapter{精确几何进阶}

\section{整数最近点对与 Minkowski 和}\label{classic-GeometryExtra}\index{Geometry\_Extra}

算法约束见紧凑版本第~\pageref{compact-GeometryExtra}~页。本节代码独立可抄；采用下列实际接口名。

\lstinputlisting[firstline=4,lastline=103]{../src/classic/geometry_extra.hpp}

\Needspace{420pt}

\section{整数三维共线、共面与方向}\label{classic-IntegerGeometry3D}\index{Integer\_Geometry\_3D}

算法约束见紧凑版本第~\pageref{compact-IntegerGeometry3D}~页。本节代码独立可抄；采用下列实际接口名。

\lstinputlisting[firstline=104,lastline=142]{../src/classic/geometry_extra.hpp}
